\documentclass[11pt,a4paper]{letter}
\usepackage{fontspec}
\setmainfont{Latin Modern Roman}
\usepackage[margin=2.5cm]{geometry}
\usepackage[hidelinks]{hyperref}
\signature{Azer Valizada}
\address{Azer Valizada\\ Technical University of Munich\\ TUM School of Engineering and Design\\ 85748 Garching bei M\"unchen, Germany\\ azer.valizada@tum.de}
\date{\today}
\begin{document}
\begin{letter}{The Editors\\ Aerospace Science and Technology}
\opening{Dear Editors,}

I submit the manuscript \emph{``Adaptive Augmentation of Incremental Nonlinear Dynamic Inversion under Rotor Faults: A Fair
Comparison of RLS, MRAC and $\mathcal{L}_1$ from Quadrotor to Quadplane''} for consideration as a research article in
\emph{Aerospace Science and Technology}.

Incremental nonlinear dynamic inversion (INDI) is now widely used on multirotors and hybrid VTOL aircraft. Its only model term,
the control-effectiveness matrix, is exactly what a rotor fault corrupts, and three families of adaptive augmentation have been
proposed to correct it: recursive least squares, predictor-based model-reference adaptive control and $\mathcal{L}_1$ adaptive
augmentation. Each has so far been shown on its own vehicle and tuned by its own authors. This manuscript compares them under one
protocol, with scenarios, metrics and a tuning rule fixed before any adaptive method was run, on a quadrotor and on a quadplane
with weighted least-squares allocation over rotors and control surfaces. The main findings are:
\begin{itemize}
\item $\mathcal{L}_1$ recovers fastest from rotor faults on both vehicles and under gusts, noise and model errors, but has the
smallest delay margin; a one-axis linear model explains why.
\item RLS reaches oracle-level control only if its measurement includes the wing's known contribution and its estimate is
bounded away from zero; even then it does not identify the fault. MRAC depends on the aerodynamic model and on air data.
\item A complete lift-rotor loss on the quadplane is recoverable only from 6--8\,m/s, a limit that a controller-independent trim
analysis shows to be physical under pitch and heading hold; continuing the transition lets the aircraft outrun a loss at lower
speed.
\end{itemize}
We believe these results are useful to readers of the journal who design fault-tolerant flight control for VTOL aircraft, and
that the fixed protocol, the reported failures of our own first interpretations and the open code make the comparison easy to
check and extend.

The work is simulation-based. All code, scenario definitions and result data are openly available at
\url{https://github.com/azervalizada/adaptive-indi-vtol}. The use of the AI assistant Claude (Anthropic) in developing the code
and in preparing the manuscript is described in the Methods section and in the AI declaration, in line with the journal's policy.
The manuscript is original, has not been published, and is not under consideration elsewhere. The author declares no competing
interests.

\closing{Yours sincerely,}
\end{letter}
\end{document}
